\documentclass{article}
\usepackage[T1]{fontenc}
\usepackage{lmodern}
\usepackage{graphicx}
\usepackage{amsmath,amssymb}
\usepackage[a4paper,margin=2cm]{geometry}
\usepackage[absolute,overlay]{textpos}
\setlength{\TPHorizModule}{1mm}
\setlength{\TPVertModule}{1mm}
\usepackage[indonesian]{babel}
\usepackage{hyperref}
\setlength{\emergencystretch}{2em}
% unit: Halaman 92

\begin{document}

% === Halaman 92 (llm) ===

\section*{Catatan:}

\begin{itemize}
    \item Untuk memutar file audio, jika belum terpasang, install terlebih dahulu
    \vspace{5mm}
    \begin{verbatim}
pip install simpleaudio
    \end{verbatim}
    \item Fungsi untuk memutar file audio:
    \vspace{5mm}
    \begin{verbatim}
import simpleaudio as sa

def play_audio(file_path):
    print(f"Memutar audio: {file_path}")
    wave_obj = sa.WaveObject.from_wave_file(file_path)
    play_obj = wave_obj.play()
    play_obj.wait_done() # Tunggu sampai selesai diputar
    \end{verbatim}
\end{itemize}

\end{document}
